\section{Conclusion}

This paper studies whether Chinese household inflation expectations exhibit diagnostic overreaction under uncertainty and whether that behavior survives identification-focused scrutiny. The evidence combines salience-IV estimation, dynamic local projections, mechanism competition against information rigidity, and Bayesian diagnostic modules. Across these layers, revisions predict subsequent forecast errors with the diagnostic sign, and the relation strengthens in high-salience or high-uncertainty states.

Policy relevance comes from computability rather than narrative interpretation. The estimated coefficients map into a real-time forecast-adjustment rule that improves backtest accuracy and interval coverage relative to naive survey forecasts, with gains reported in explicit RMSE and MAE units. That mapping gives central-bank communication a measurable risk metric: when salience spikes, unadjusted expectations are more likely to overstate near-term inflation.

The inference remains sample-bound. The design uses national quarterly data from 2011Q1 to 2025Q3, lacks micro household panels for direct exposure heterogeneity, and relies on external salience proxies that are plausible but not experimentally assigned. Future work can extend the same identification logic to province-household matched data and richer media microstructure, where staggered exposure designs and stronger exclusion tests can be implemented without changing the mechanism definitions used here.
